% WLC: exact finite-place Weil distribution, including its unit contribution.
\section{Every finite-place Weil term and the conductor operator}
\label{sec:cc-local-weil-conductor}

\paragraph{Original sources and scope.}
The local distribution and its multiplicative Haar measure are from
Alain Connes, \href{https://arxiv.org/abs/math/9811068v1}{\emph{Trace
formula in noncommutative geometry and the zeros of the Riemann zeta
function}}, original author TeX, local trace theorem (12)--(17),
lines 1367--1436; VIII(10), lines 2353--2360; and Appendix II(9),
(22)--(29), lines 4280--4495. The original Weil principal value and
its conductor term are human results, not new discoveries here.
The finite-field computations and the receiving operator below
are proved explicitly with all measures retained.
Alain Connes, Caterina Consani and Matilde Marcolli,
\href{https://arxiv.org/abs/math/0703392v1}{\emph{The Weil proof and
the geometry of the adeles class space}}, original author locators
\texttt{Weilexpl}, \texttt{princvalrhobeta}, \texttt{TraceformulaH1}
and \texttt{intalphae}, give the global receiving formula and its
differential-idele term. That original trace theorem is a cited
input; the calculation here evaluates its finite local terms.

\subsection{The original measure and an explicit principal value}

Let $F$ be a nonarchimedean local field, $\mathcal O$ its ring of
integers, $\varpi$ a uniformizer, and $q=|\mathcal O/\varpi\mathcal O|$.
Retain additive Haar measure $dx$ with $\operatorname{vol}(\mathcal O)=1$
and $|\varpi|=q^{-1}$. The source's multiplicative measure is
\begin{equation}\tag{WLC1}
 d^*x=\frac{\log q}{1-q^{-1}}\frac{dx}{|x|},\qquad
 \operatorname{vol}_{d^*x}(\mathcal O^\times)=\log q.
\end{equation}
Indeed every shell $\varpi^j\mathcal O^\times$ has that same
multiplicative mass. The integral on $1\le |x|\le q^M$ is
$(M+1)\log q$, whose ratio to $\log(q^M)$ tends to one,
as required by the original measure convention. No unit mass
has been substituted for the displayed $\log q$.

For a locally constant compactly supported function $b$ on $F^\times$
define the original Weil principal value by the explicit formula
\begin{equation}\tag{WLC2}
\begin{split}
 W_F(b)={}&\int_{|u|\ne1}\frac{b(u^{-1})}{|1-u|}\,d^*u\\
 &+\int_{\mathcal O^\times}
       \frac{b(u^{-1})-b(1)}{|1-u|}\,d^*u.
\end{split}
\end{equation}
Both integrals are finite: only finitely many valuation shells
meet the first support, and the second numerator vanishes on a
neighborhood of one. This agrees with ordinary integration away
from one and assigns zero to the test $1_{\mathcal O^\times}$.
It is the source's Appendix II(9) principal value. Its equality
with the additive-character prescription is the cited Lemma 2
in that appendix, for a character trivial on $\mathcal O$ but
nontrivial on $\varpi^{-1}\mathcal O$.

The last condition is the conductor condition. It does not assert
that the additive character's entire kernel is $\mathcal O$.
The wording $\ker e_F=\mathcal O$ in the later CCM passage is
too strong for general local fields. For example, on
$\mathbb F_p((T))$, the character
$\beta(\sum a_jT^j)=\exp(2\pi i a_{-1}/p)$ has the stated
conductor but $\beta(T^{-2})=1$. In fact every additive character
of this field takes values in the $p$th roots of unity, since
the additive group has exponent $p$; its quotient by its kernel
has at most $p$ elements. The quotient
$\mathbb F_p((T))/\mathbb F_p[[T]]$ is infinite, so no such
character can have that ring as its entire kernel. WLC2 uses
the valid original conductor prescription, not that stronger
assertion.

The exact Fourier relation survives with that conductor condition.
One has $\mathcal O^\perp=\mathcal O$: an element of $\mathcal O$
pairs trivially with every element of $\mathcal O$. If
$v(x)=-k<0$, then $x\mathcal O=\varpi^{-k}\mathcal O$ contains
$\varpi^{-1}\mathcal O$, where the character is nontrivial;
hence some $y\in\mathcal O$ has $\beta_0(xy)\ne1$.
Integration of this nontrivial character over $\mathcal O$
is zero by translation invariance. Therefore the actual Fourier
map satisfies
\[
 \int_{\mathcal O}\beta_0(xy)\,dy=1_{\mathcal O}(x).
\]
This is the comparison used in Connes's Appendix II(24), with
the original additive measure of mass one on $\mathcal O$.

Let $\chi:F^\times\to\mathbb C^\times$ be a continuous character
whose restriction to $\mathcal O^\times$ has finite image.
Put $\lambda=\chi(\varpi)\ne0$, and let its unit conductor be
\begin{equation}\tag{WLC3}
 a=0\ \text{if }\chi|_{\mathcal O^\times}=1;
 \qquad
 a=\min\{j\ge1:\chi|_{1+\varpi^j\mathcal O}=1\}
 \ \text{otherwise}.
\end{equation}
For a continuous unitary character the finite-image condition
follows from continuity and the profinite unit group: an open
subgroup maps into a sufficiently short arc with no nontrivial
subgroup of the unit circle, and therefore into one.
For $h\in C_c^\infty(\mathbb R_{>0})$ use the actual test
\begin{equation}\tag{WLC4}
 b(x)=h(|x|)\chi(x).
\end{equation}
It is locally constant on $F^\times$ and compactly supported there.

\subsection{Every shell and the ramified unit term}

On $u=\varpi^j\epsilon$, $j>0$, one has
$|1-u|=1$ and
$b(u^{-1})=h(q^j)\lambda^{-j}\chi(\epsilon)^{-1}$.
For $u=\varpi^{-j}\epsilon$ the denominator is $q^j$ and
the numerator is $h(q^{-j})\lambda^j\chi(\epsilon)^{-1}$.
The integral of a nontrivial unit character is zero: translating
the Haar integral by a unit with character value different from
one multiplies it by that value without changing the integral.
For $a=0$, the unit numerator in WLC2 vanishes identically.
Consequently
\begin{equation}\tag{WLC5}
 W_F(b)=\log q\sum_{j\ge1}
 \left[\lambda^{-j}h(q^j)+q^{-j}\lambda^jh(q^{-j})\right]
 \qquad(a=0).
\end{equation}
The sum has finitely many nonzero terms for the stated $h$.

For $a>0$, every off-unit shell instead integrates to zero.
Write $K_j=1+\varpi^j\mathcal O$, $j\ge1$. Their original masses are
\begin{equation}\tag{WLC6}
 \mu^*(K_j)=\frac{\log q}{1-q^{-1}}q^{-j},\qquad
 \int_{K_j}\chi(u)^{-1}d^*u=
 \begin{cases}0,&j<a,\\\mu^*(K_j),&j\ge a.\end{cases}
\end{equation}
The mass follows from additive translation by one and $|u|=1$.
The character integral follows from the same translation argument,
now in the compact subgroup $K_j$, and from WLC3.
For $u\ne1$ in the unit group the exact finite-at-each-point identity is
\begin{equation}\tag{WLC7}
 \frac1{|1-u|}=1+
       \sum_{j\ge1}(q^j-q^{j-1})1_{K_j}(u).
\end{equation}
If the valuation of $1-u$ is $k$, the sum stops at $k$ and
telescopes to $q^k$. After multiplication by
$\chi(u)^{-1}-1$, every term with $j\ge a$ is zero, so
the following integration is a finite calculation:
\begin{equation}\tag{WLC8}
\begin{split}
 \int_{\mathcal O^\times}\frac{\chi(u)^{-1}-1}{|1-u|}d^*u
 &=-\log q-
  \sum_{j=1}^{a-1}(q^j-q^{j-1})
                    \frac{\log q}{1-q^{-1}}q^{-j}\\
 &=-a\log q.
\end{split}
\end{equation}
Thus the complete ramified value, including the singular unit stratum, is
\begin{equation}\tag{WLC9}
 W_F(b)=-a\log q\,h(1)\qquad(a>0).
\end{equation}
Its local Euler factor is one, but its local Weil distribution
has not vanished. WLC8 independently recovers the human conductor
formula Connes VIII(10), retaining its original sign and measure.

If the additive character is $\beta_c(x)=\beta_0(cx)$ for a
conductor-$\mathcal O$ reference character $\beta_0$, the cited
CCM identity \texttt{intalphae} retains the additional term
\begin{equation}\tag{WLC10}
 W_{F,\beta_c}(b)=W_{F,\beta_0}(b)+\log|c|\,b(1).
\end{equation}
No different-idele term is removed by the shell computation.

\subsection{The original rational prime and Dirichlet dictionaries}

Use $C_{\mathbb Q}=\mathbb R_{>0}\times\prod_r\mathbb Z_r^\times$
in the exact PLP20 coordinates. A finite character $\chi$ of
$(\mathbb Z/N\mathbb Z)^\times$ induces a compact-unit character
$\widetilde\chi$ and a unique primitive character $\chi_*$
of conductor $f_\chi\mid N$. For
\begin{equation}\tag{WLC11}
 f(y,u)=h(y)\widetilde\chi(u),\qquad
 b_p(x)=f(\iota_p(x)),
\end{equation}
the local uniformizer has the exact coordinates
\begin{equation}\tag{WLC12}
 \iota_p(p^j\epsilon)=
 \left(p^{-j},\ (u_p=\epsilon,\ u_r=p^{-j}\ (r\ne p))\right).
\end{equation}
It follows that the local unit conductor is
$a_p=v_p(f_\chi)$. If $a_p=0$, then
$\widetilde\chi(\iota_p(p))=\chi_*(p)^{-1}$.
This follows by multiplying the finitely many other unit components;
their integer evaluation at $p$ is exactly $\chi_*(p)$.
The inverse local character is essential: WLC5 therefore becomes
\begin{equation}\tag{WLC13}
 W_p(b_p)=\log p\sum_{j\ge1}
 \left[\chi_*(p)^jh(p^j)+p^{-j}\chi_*(p)^{-j}h(p^{-j})\right]
 \quad(p\nmid f_\chi),
\end{equation}
whereas at a conductor prime its full value is
\begin{equation}\tag{WLC14}
 W_p(b_p)=-v_p(f_\chi)\log p\,h(1).
\end{equation}
This agrees with the two point-return and inverse-pullback rows
PHS19--21, rather than identifying their opposite phases.

Summing the finite places gives the exact expression
\begin{equation}\tag{WLC15}
\begin{split}
 \sum_pW_p(b_p)={}&
 \sum_{p\nmid f_\chi}\log p\sum_{j\ge1}
 \left[\chi_*(p)^jh(p^j)+p^{-j}\chi_*(p)^{-j}h(p^{-j})\right]\\
 &-\sum_{p\mid f_\chi}v_p(f_\chi)\log p\,h(1).
\end{split}
\end{equation}
All prime-power sums here are finite: the compact support of $h$
bounds both $p^j$ and $p^{-j}$ away from their respective escaping
endpoints. There are finitely many conductor primes. The last
sum equals $-\log(f_\chi)h(1)$, with all its local summands
displayed above.

The compact character depends on its conductor, not on an extra
imprimitive modulus. For $p\mid N$ but $p\nmid f_\chi$, WLC13
still occurs. This is compatible with the original Dirichlet series
only after retaining every missing-prime factor:
\begin{equation}\tag{WLC16}
 L(s,\chi)=L(s,\chi_*)
 \prod_{\substack{p\mid N\\p\nmid f_\chi}}
          (1-\chi_*(p)p^{-s}),\qquad \Re s>1.
\end{equation}
The proof is term-by-term Euler-factor comparison in the region
of absolute convergence, as proved in PHS13--18. Replacing
$\chi_*(p)$ by the imprimitive value zero in WLC13 would lose
these actual arithmetic terms.

For the standard global additive character of $\mathbb A_{\mathbb Q}/
\mathbb Q$, the finite reference conductors are $\mathbb Z_p$.
The original CCM global trace receives WLC15 in the formula
\begin{equation}\tag{WLC17}
\begin{split}
 \operatorname{Tr}(\underline\vartheta_m(f)|H^1)
 ={}&\widehat f(0)+\widehat f(1)-\log|\mathfrak a|\,f(1)\\
 &-W_{\mathbb R,e_{\mathbb R}}(f\circ\iota_\infty)
 -\sum_p W_p(f\circ\iota_p).
\end{split}
\end{equation}
Here $\mathfrak a$ is the source's differential idele, and the
archimedean term is its specified original principal value, not
an omitted remainder. With $d^\times y=dy/y$ and compact-unit
Haar mass one, the endpoint integrals for WLC11 are exactly
\begin{equation}\tag{WLC18}
 \widehat f(0)=\delta_{\chi_*,1}\int_0^\infty h(y)\frac{dy}{y},\qquad
 \widehat f(1)=\delta_{\chi_*,1}\int_0^\infty h(y)\,dy,
 \qquad f(1)=h(1).
\end{equation}
These follow from the integral of the compact character. The
archimedean restriction is
$f(\iota_\infty(x))=h(|x|)\chi_*(\operatorname{sign}x)$;
its parity has not been discarded. For the displayed rational
reference choice $|\mathfrak a|=1$, but WLC17 retains its
original location and WLC10 its local dependence. WLC17 invokes
the cited global theorem; WLC1--16 evaluate precisely its finite
part and do not replace that theorem's analytic proof.

\subsection{An exact finite conductor operator before the trace}

On $V_N=\mathbb C^{U_N}$ use the original finite Haar measure
$\mu_N(u)=1/\varphi(N)$, and the PHS28 Fourier coefficients.
For every prime divisor $p$ of $N$, put $r=v_p(N)$ and define
\begin{equation}\tag{WLC19}
 C_{p,N}\chi=v_p(f_\chi)\log p\,\chi,\qquad
 C_N=\sum_{p\mid N}C_{p,N}.
\end{equation}
This defines genuine linear operators by the finite character
basis. Their full action and quadratic form are
\begin{equation}\tag{WLC20}
\begin{gathered}
 (C_Na)(u)=\sum_{\chi\in\widehat U_N}
       \left(\sum_{p\mid N}v_p(f_\chi)\log p\right)
                       \widehat a(\chi)\chi(u),\\
 \langle a,C_Na\rangle_{\mu_N}
  =\sum_{\chi\in\widehat U_N}
       \left(\sum_{p\mid N}v_p(f_\chi)\log p\right)
                     |\widehat a(\chi)|^2\ge0.
\end{gathered}
\end{equation}
Orthogonality proves the formula and self-adjointness, with all
measure factors inherited from PHS9 and PHS28. The kernel is
the constant-character line, since conductor one means the
character is trivial. This finite positive quadratic form is
not a positivity assertion for WLC17's entire trace.

For the general finite-character test $f(y,u)=h(y)a(u)$, linearity
and WLC14 show that its full ramified-unit contribution is
\begin{equation}\tag{WLC21}
 -h(1)(C_Na)(1).
\end{equation}
Evaluation at one sums the character coefficients because every
character has value one there. Unramified valuation-shell terms
and WLC17's other contributions remain separate summands of
the same original trace.

The multiplicities can also be computed exactly. On $U_{p^r}$
there is one conductor-zero character. The number with conductor
at most $j$, $1\le j\le r$, is $\varphi(p^j)$, since reduction
onto $U_{p^j}$ is surjective and its characters are exactly the
pullbacks. Counting $a=\sum_{j=0}^{r-1}1_{a>j}$ gives
\begin{equation}\tag{WLC22}
 \operatorname{Tr}(C_{p,N})=
 \varphi(N/p^r)\log p
 \left[r\varphi(p^r)-1-\sum_{j=1}^{r-1}\varphi(p^j)\right].
\end{equation}
The factor $\varphi(N/p^r)$ counts the independent characters
at all other prime factors. The displayed finite sum includes
the empty-sum case $r=1$, so $p=2,r=1$ correctly contributes
zero. Telescoping $\varphi(p^j)=p^j-p^{j-1}$ makes the bracket
$r\varphi(p^r)-p^{r-1}$; both the local degree and every
filtration contribution were retained in WLC22.

For the entire original bounded distributive lattice, the maps are
\begin{equation}\tag{WLC23}
 C_N^\uparrow:G_L(V_N)\longrightarrow G_L(V_N),\quad
 (a,\lambda)\longmapsto(C_Na,\lambda),\qquad
 W_F^\uparrow(b,\lambda)=(W_F(b),\lambda).
\end{equation}
Linearity proves additivity and compatibility with
$G_L(\mathbb C)$ scalars. A nonzero output amplitude comes
from a nonzero input, hence has its forced top label; a zero
output keeps its original label. These are common-label
semimodule maps, not arbitrary basis changes on independently
labelled coordinate products. In particular the constant test
with top label satisfies
\begin{equation}\tag{WLC24}
 C_N^\uparrow(1,1_L)=e,\qquad e\ne\tau
 \quad\text{for nontrivial }L.
\end{equation}
At every support stalk, its map commutes with the original
localization $(a,\lambda)\mapsto[\lambda]_J$ because it keeps
that label. No measure on the arbitrary lattice has been
invented. The unit conductor, prime-power phases, omitted-prime
factors and labelled zero images have all been evaluated in
their specified original objects.
